\documentclass{beamer}
\usetheme{Singapore}
\usecolortheme{orchid}

% Modern Purple Color Scheme
\definecolor{deepPurple}{RGB}{75, 0, 130}
\definecolor{vibrantPurple}{RGB}{138, 43, 226}
\definecolor{lightPurple}{RGB}{200, 150, 255}

\setbeamercolor{primary}{bg=deepPurple, fg=white}
\setbeamercolor{frametitle}{bg=vibrantPurple, fg=white}
\setbeamercolor{normal text}{fg=deepPurple}
\setbeamercolor{structure}{fg=vibrantPurple}
\setbeamercolor{alerted text}{fg=vibrantPurple}

% Fun Modern Font
\usepackage[sfdefault]{carlito}
\renewcommand{\familydefault}{\sfdefault}

\title{Applications of Market Design}
\subtitle{Chapter 4}
\author{Alyssa Alvarez}
\date{\today}

\begin{document}

\frame{\titlepage}

\begin{frame}{What is Market Design?}
\textbf{Goals of market design research:}
\begin{enumerate}
  \item Identify institutional sources of inefficiency
  \item Evaluate alternative market rules
  \item Propose and implement improved designs
\end{enumerate}
\vspace{1cm}
Market design applies \textit{theoretical} and \textit{empirical} tools to real markets to improve outcomes.
\end{frame}

\begin{frame}{Diagnosing Market Failures}
\textbf{Key insight:} Market failures can arise from \textbf{design flaws}, not just market power.
\vspace{0.5cm}
\\
\textbf{Common diagnostic focus:}
\begin{itemize}
  \item How information is aggregated
  \item How agents are allowed to act
  \item How incentives are shaped by institutional rules
\end{itemize}
\vspace{0.5cm}
\textbf{Diagnosis requires:}
\begin{itemize}
  \item Close examination of clearing rules
  \item Empirical analysis to distinguish causes
\end{itemize}
\end{frame}

\begin{frame}{Coordination Failures and Centralized Clearing}
\textbf{Problems with decentralized markets:}
\begin{itemize}
  \item Congestion and mismatches
  \item Market unraveling (early contracting to avoid uncertainty)
  \item Participants remain unmatched despite gains from trade
\end{itemize}
\vspace{0.5cm}
\textbf{Benefits of centralized clearing:}
\begin{itemize}
  \item Aggregates dispersed preference information
  \item Reduces uncertainty about outcomes
  \item Improves match quality
\end{itemize}
\end{frame}

\begin{frame}{Strategic Behavior and Incentive Problems}
\textbf{Core challenge:} Poorly designed mechanisms create incentives for \textbf{strategic misrepresentation}.
\vspace{0.5cm}
\\
\textbf{Consequences:}
\begin{itemize}
  \item Reduced efficiency
  \item Harms less informed or sophisticated participants
\end{itemize}
\vspace{0.5cm}
\textbf{Design principle:} Aim for \textbf{strategy-proofness}, where truthful revelation is optimal.
\vspace{0.5cm}
\\
\textbf{Behavioral realism:} Mechanisms must account for mistakes, not just theoretical equilibrium.
\end{frame}

\begin{frame}{Evaluating and Comparing Market Designs}
\textbf{Central task:} Compare alternative institutional rules
\vspace{0.5cm}
\\
\textbf{Evaluation criteria:}
\begin{itemize}
  \item Allocative efficiency
  \item Distributional outcomes
  \item Incentive compatibility
  \item Participation and access
\end{itemize}
\vspace{0.5cm}
\textbf{Methods:}
\begin{itemize}
  \item Structural empirical models
  \item Counterfactual simulations
  \item Welfare analysis
\end{itemize}
\end{frame}

\begin{frame}{Efficiency vs. Fairness Tradeoffs}
\textbf{Challenge:} Multiple normative objectives often conflict
\vspace{0.5cm}
\\
\textbf{Fairness considerations:}
\begin{itemize}
  \item Respecting priority rights
  \item Avoiding justified envy
  \item Ensuring equitable access
\end{itemize}
\vspace{0.5cm}
\textbf{Key insight:} Institutional rules encode \textbf{social and political values}, not just efficiency goals.
\vspace{0.5cm}
\\
Market design provides a framework for explicitly analyzing these tradeoffs.
\end{frame}

\begin{frame}{Designing New Market Mechanisms}
\textbf{When are new designs needed?}
\begin{itemize}
  \item Constraints are complex
  \item Preferences are over bundles
  \item Transfers are restricted or prohibited
\end{itemize}
\vspace{0.5cm}
\textbf{Requirements for successful mechanisms:}
\begin{itemize}
  \item Respect incentive constraints
  \item Be computationally tractable
  \item Be understandable and usable by participants
\end{itemize}
\vspace{0.5cm}
\textbf{Examples:} Kidney exchange, course allocation, humanitarian resource allocation
\end{frame}

\begin{frame}{Complementarity of Theory and Empirics}
\textbf{Theory:}
\begin{itemize}
  \item Identifies feasible allocations
  \item Clarifies incentive and informational constraints
\end{itemize}
\vspace{0.5cm}
\textbf{Empirics:}
\begin{itemize}
  \item Measures welfare impacts
  \item Tests behavioral assumptions
  \item Guides policy decisions
\end{itemize}
\vspace{0.5cm}
\textbf{Best practice:} Combine mechanism design theory, structural estimation, and field data.
\end{frame}

\begin{frame}{Institutional Adoption and Political Economy}
\textbf{Challenge:} Well-designed mechanisms may not be adopted.
\vspace{0.5cm}
\\
\textbf{Sources of resistance:}
\begin{itemize}
  \item Incumbent institutions
  \item Participants who benefit from inefficiencies
\end{itemize}
\vspace{0.5cm}
\textbf{Implications:} Market design intersects with political economy, regulation, and governance.
\vspace{0.5cm}
\\
\textbf{Reality:} Government intervention may be necessary to implement socially superior designs.
\end{frame}

\begin{frame}{Core Lessons from Chapter 4}
\begin{enumerate}
  \item \textbf{Institutions shape outcomes} --- Design matters as much as competition
  \item \textbf{Diagnose carefully} --- Understand the true source of market failures
  \item \textbf{Evaluate systematically} --- Use theory and data to compare alternatives
  \item \textbf{Balance objectives} --- Efficiency and fairness are not always compatible
  \item \textbf{Design with constraints} --- Account for incentives, computation, and usability
  \item \textbf{Combine tools} --- Theory and empirics are complementary
  \item \textbf{Plan for adoption} --- Political economy matters for implementation
\end{enumerate}
\end{frame}

\end{document}
